\section{相关工作}\label{sec:related}
\subsection{不完整观测下的重建与预测}\label{sec:2_1}

不完整观测建模包含历史恢复与未来预测两个相关但不同的目标。GRIN 通过图循环计算和消息传递重建多变量序列，PriSTI 则利用条件扩散和时空关联建模缺失数据~\citep{cini2022filling,liu2023pristi}。二者说明，关联位置能够为局部缺失提供信息来源；但重建历史的目标与预测未来并不相同，插补质量也不能直接代替预测性能评价。

面向未来预测，GCN-M 联合建模缺失处理与交通预测，通过观测掩码、多尺度历史记忆和图卷积利用局部与全局信息~\citep{zuo2023graph}。GinAR 针对部分变量历史不可用的情况，将插值注意力和自适应图卷积引入递归预测~\citep{yu2024ginar}。UIGNN 同时面向缺失位置预测与不确定性估计，与未观测位置泛化形成交叉~\citep{mei2023uncertainty}。这些工作已经建立了缺失处理和空间推理的联系，并以掩码区分可见与缺失读数。本文在此基础上关注 OS 的时间结构：可见读数数量相同但位于窗口内不同位置时，预测所依据的信息并不相同；重建之外，观测角色与状态进度如何参与传播组合，跨节点上下文如何进一步条件化未来预测。相应地，实验分别干预状态进度、OS 条件化传播与跨节点上下文，以区分各机制的贡献，并在固定可见数量下比较不同时间布局。

\subsection{归纳与部分感知预测}\label{sec:2_2}

归纳建模研究如何将已观测位置上学习的规律用于缺少训练信号的位置。IGNNK 通过子图采样学习可迁移的消息传递，解决未采样位置的时空插值~\citep{wu2021inductive}；UIGNN 和 MoGERNN 则进一步预测未来交通状态~\citep{mei2023uncertainty,zhou2025mogernn}。其中，MoGERNN 采用图专家混合与图循环编码—解码结构，并考察传感器加入和失效情形，说明变化的传感器配置已是归纳交通预测的重要研究对象。MoGERNN 依据节点输入特征稀疏地选择图专家，并在进入编码—解码器之前为未观测节点构造表示；ALFGraph 则把传播算子的组合放进图循环单元的每个门，由逐步的观测角色、状态进度与掩码稠密路由，不读取信号值。

不同方法对训练和推理阶段的观测条件有不同假设。SLPF 研究长期部分感知预测，利用训练期较充分的数据学习向推理期未感知位置的空间转移~\citep{liu2025spatiotemporal}。SIML-MSIGN 将不同观测布局组织为任务，通过掩码门控扩散与结构归纳元学习开展跨布局预测及适应~\citep{zhao2026structure}。SC-SAMG 利用少量目标位置数据生成节点与图表示，并进行少步参数适应~\citep{wang2026cold}；其中的“支持”指用于适应的目标样本，与本文窗口内的 OS 含义不同。因此，本文并不将“变化观测下的预测”视为未经研究的领域。ALFNet 的侧重点是固定参数下的窗口级 OS 条件化：在统一输入预算内，研究可用历史的状态和布局如何影响空间—时间计算。总体评价考察这一设计在归纳设置中的效果，已知节点配对实验则用于分离 OS 变化本身带来的预测响应。

\subsection{空间依赖与时空表示学习}\label{sec:2_3}

空间依赖建模为交通预测提供了利用关联位置的基本途径。STGCN 结合图卷积与时间卷积；Graph WaveNet 和 AGCRN 利用可学习表示构建自适应关系，降低对固定邻接的依赖~\citep{yu2018spatio,wu2019graph,bai2020adaptive}。PDFormer 将传播时延纳入长程空间交互，PatchSTG 通过空间分块及块内、块间交互组织局部与全局信息~\citep{jiang2023pdformer,fang2025efficient}。这些方法分别改善空间关系、交互范围和计算粒度，但学习到自适应图，并不自动意味着传播会根据每个窗口的 OS 进行调整。空间交互还需与时间表示能力共同考察。STID 通过简洁的时空身份表示建立有竞争力的预测基线~\citep{shao2022spatial}；STBP 利用可扩展上下文模式库支持城市持续预测~\citep{liu2026general}。本文关注的不是无限扩大空间交互范围或在线更新模式库，而是如何将可用的跨节点历史用于当前预测条件化。ALFTemporal 生成的上下文经 FiLM~\citep{perez2018film} 调制图循环预测，ALFGraph 则承担以 OS 为条件的空间传播。实验中的同容量局部上下文替代直接检验跨节点信息的作用，避免将模型收益仅归结为更大的表示规模。
